% 论语学习与解读.tex
% Auto-generated from 感悟论语.odt by tools/odt_to_tex.py

\setpartplant{artwork/tangdi.png}
  {唐棣之华 · 《论语·子罕》}
  {未之思也，夫何远之有？}
  {《论语·子罕》}
  {论语学习与解读}
\part{论语学习与解读}

\chapter{论语入门读书顺序}

1. 傅佩荣《论语300讲》

2. 杨伯峻《论语译注》

3. 钱穆《论语新解》

4. 李泽厚《论语今读》

5. 朱熹《四书章句集注》论语部分

6. 傅佩荣《孔子辞典》

其中，傅佩荣《孔子辞典》与杨伯峻《论语译注》作为辞典随时查询。

\entrydate{2017.02.19}

\chapter{由《论语300讲》所感}

听过与读过傅佩荣教授在《论语300讲》中的解读后，我对《论语》有了更为真切的体会：不只是在自我管理与为人处事方面获益匪浅，更是促成了对人生方向的探索并使其逐渐明晰，以及对此方向锲而不舍的坚守与追寻。从孔子与其弟子的对话可以看出，古时读书人在求学伊始的第一要务是立志，是以修身、做人为本，而后才是求知——君子务本，本立而道生。由于他们能够不论在何时何地所作的抉择都是围绕着这个志向，因而较之今天的年轻人也就少有人生、事业与感情上的困惑。这种“择善固执”无需顾及外在环境与客观条件的优劣——饭疏食饮水，曲肱而枕，乐亦在其中矣；一箪食，一瓢饮，在陋巷，回也不改其乐，亦不会受制于时势命运——君子无终食之间违仁，造次必于是，颠沛必于是。而且，能够以毕生的付出和牺牲坚持走在通往至善与成为君子的道路上，直至最终“止于至善”。或者退一步说，即便这样的终极目标因天命而无法实现，然而正因为“朝闻道，夕可死矣”，所以纵然在刚刚意识到人生正途何在的时候便被上天剥夺了实现的机会，也同样是为时不晚、幸运且不留遗憾的。

以上便是傅佩荣在《论语300讲》中反复予以解读和强调的主旨思想：“人性向善，择善固执，止于至善”。这是值得我们每个中国人用心体会、慎思笃行和贯彻一生的精神内核。

\begin{classic}
《论语·学而篇第一 1.2》 有子曰：「其为人也孝弟，而好犯上者，鲜矣；不好犯上，而好作乱者，未之有也。君子务本，本立而道生。孝弟也者，其为仁之本与！」

《论语·卫灵公篇第十五 15.2》 在陈绝粮。从者病，莫能兴。子路愠见曰：「君子亦有穷乎？」子曰：「君子固穷，小人穷斯滥矣。」

《论语·述而篇第七 7.16》 子曰：「饭疏食，饮水，曲肱而枕之，乐亦在其中矣。不义而富且贵，于我如浮云。」

《论语·雍也篇第六 6.11》 子曰：「贤哉，回也！一箪食，一瓢饮，在陋巷，人不堪其忧，回也不改其乐。贤哉，回也！」

《论语·里仁篇第四 4.5》 子曰：「富与贵，是人之所欲也，不以其道得之，不处也。贫与贱，是人之恶也，不以其道得之，不去也。君子去仁，恶乎成名。君子无终食之间违仁，造次必于是，颠沛必于是。」

《论语·里仁篇第四 4.8》 子曰：「朝闻道，夕死可矣！」
\end{classic}

\entrydate{2018.02.06}

\chapter{《论语》的解读}

\begin{classic}
《论语·学而篇第一 1.1》 子曰：「学而时习之，不亦说乎？有朋自远方来，不亦乐乎？人不知而不愠，不亦君子乎？」
\end{classic}

以上《论语 学而》第一章大家已是耳熟能详了。记得这一章连同“学而不思则罔，思而不学则殆”等曾经被选入我初中一年级的语文课本。虽然这一章看似浅显易懂，但是否能够理解到位乃至切实践行，就不好说了。仅以“学而时习”中的“时”字为例，一般人望文生义，会将其理解为“时常”，也就像是中小学生的定期复习那样，要想学得好，就要不断反复。这样的理解对于以其四分之一至三分之一人生投入应试教育的广大学生无疑有着“重大”的教益，因此也颇得老师与家长的认可。

相比之下，大师的解读便没有那么功利了。例如，钱穆在《论语新解》中对“时”字给出了三种解释：

\begin{itemize}
\item 指年岁。例如“古人六岁始学识字，七八岁教以日常简单礼节，十岁教书写计算，十三岁教诗歌舞蹈”。
\item 指季节。古人春夏学诗乐弦歌，秋冬学书礼射猎。
\item 指一日中的晨夕。
\end{itemize}

可以看出，钱穆的解释采用的是“时”字的字面意思，即“时间”，只不过尺度不同罢了。看过这样的解释，便会对古人学习的内容、流程、习惯有所了解，相当于掌握了时代背景。

而傅佩荣教授在其《论语三百讲》中则将“时”解释为“恰当的时候”。意即对学所学的礼仪、射猎等知识和技术、技能，需要在合适的时候加以练习。这是因为，孔子所传授的内容并非可以独立于现实的纯粹理论和逻辑推理，而是需要在生活中加以实践和反省的人生学问与哲学，例如人生志向的确立、原则的坚守、孝敬的践行。它们都是依赖于环境与客观条件的，甚至取决于一定的机缘。

由以上可以看出，对于《论语》的学习与理解，普通人与大师有所不同，大师与大师之间也不尽相同，正所谓仁者见仁、智者见智。对此，我认为不必拘泥和固执于所谓的权威、原意、真意，而应以能够对书中的词汇、内容加以发挥、引申，从而触发个人感悟并由感悟落实到积极行动为好。若从负面的例子来看，经过几千年的演变，如今人们头脑中的儒家思想早已不是原本的样子，尤其是掺入了许多僵化、教条甚至是腐朽的观念，其影响还流传得相当广。但事实上，这些错误、片面的理解根本不是孔子最初的思想。这样看来，《论语》在事实上已经成为了囊括不同理解的统一形式。为此，为什么不能用提炼数学理论的方式来看待、处理《论语》？即，将其当作由抽象性符号构建的框架，用以容纳来自不同视角的解读、阐释。同时，要求这些阐释必须具有积极与正面的意义，深刻而丰富内涵，既可以从宏观把握、又可以指导细处的实践。由此，《论语》将成为不断完善、深化的思想体系，无论对于个人的修身、立志，还是对于优秀民族性格的构建，都是大有裨益的。

\entrydate{2017.07.08}

\chapter{《论语·学而1.1》解读}

《论语》的编排以“学而”作为首篇，以“学而时习”开宗明义，可以看出儒家对于“学习”的重视。其中，“时”是指“在适当的时候”（杨伯峻《论语译注》 P1）或“时常”、“时时”（朱熹《四书章句集注》 P47，谢氏曰：「无时而不习」）。傅佩荣(《论语300讲》上册 P1）则进一步强调学习要把握时机，在适当的时候去实践。当然，这里的学习不仅指科学知识，还应拓展到生活、社会、待人、做事等诸多方面，涉及生命、人生的大学问、大智慧。它们彼此相辅相成，是贯穿一个人毕生追求君子之道、推己及人、止于至善的基石。

具体说来：

\begin{itemize}
\item 在学习与研究方面，要做好详细的计划，这样才能做到有规律的时常复习、练习，而不是浅尝辄止、一曝十寒。同时，要注意不能避重就轻，光挑那些相对容易的科目，例如，偏向实用而忽略理论；或是见异思迁，只选那些愉悦感观的材料，例如诗词歌赋而非思辩哲学。随着年龄的增长，我们积累的知识与技能越来越多，再向其中系统性地增添大量新内容也会愈加困难。此时要想做到“月无忘其所能”，就必须增加复习的次数与频率。
\item 在行仁与孝敬方面，需要把握恰当的时机。在每一次与家人的沟通交流中，在待人接物中，在职业发展与人生选择中，珍惜现实提供的机会，回顾并应用从论语中学到的原则。
\item 在个人修行方面，要在恰当的时候，例如遇到倦怠、松懈、怯懦、灰心等个人性格弱点出现的情况时，跳出自己的舒适圈，与坏习惯做斗争。
\end{itemize}

由此推想开来：能将学习本身定位为人生的基本要素，便不会将其当作追名逐利的敲门砖。如果一个人发自内心地重视学习，其必然会坚持真理，不盲信、不盲从，拥有超越评价体系（游戏规则）的独立思考，做事亦会有原则和底线。将学习放在首位，才可以始终怀有童心、初心，以开放的胸怀了解、接纳新事物，才能做到苟日新、日日新。以学习作为终极目标，才能真正做到慎思笃行，始终走在择善固执、止于至善的道路上。能成就这样的自我，便是做到了孔子所说的“君子不器”。而且同样地，亦不会把他人当作工具。所谓“己欲立而立人，己欲达而达人”、“君子成人之美”是也。而这些均与学习是密不可分的。

有了书本上和实践中笃实的学习乃至形成深厚的积淀，志同道合的朋友相聚时，才会“以文会友，以友辅仁”。彼此交流的是关于永恒真理的理解与感悟。而那些酒桌上常见的热烈与殷勤，日常相处的如胶似漆，充其量只能算是小人的“同而不和”罢了。（注：《论语》中的小人不含贬义，只是用来形容一个人的言行、心智仍像是孩童，就像是武志红提出的“巨婴”概念）人生往往难觅投缘的好友，更难觅“爱真理胜过吾师”的知己。这样看来，“有朋自远方来”的情境岂不是一大快乐？

“学而”的最后一部分，“人不知而不愠，不亦君子乎”，讲的是宠辱不惊的人生境界。知音难觅，而误解司空见惯。能够形成这样的豁达，有赖于不断的学习和实践来给自己提供足够的智慧和修养。不仅如此，孔子亦说过“不患人之不己知，患己之不能也”，由是观之，“人不知而不愠”恰恰又是鞭策自我更进一步虚心与深入学习的动力。

\begin{classic}
《论语·学而篇第一 1.1》 子曰：「学而时习之，不亦说乎？有朋自远方来，不亦乐乎？人不知而不愠，不亦君子乎？」

《论语·为政篇第二 2.12》 子曰：「君子不器。」

《论语·雍也篇第六 6.30》 子贡曰：「如有博施于民，而能济众，何如？可谓仁乎？」子曰：「何事于仁，必也圣乎！尧舜其犹病诸！夫仁者，己欲立而立人，己欲达而达人。能近取譬，可谓仁之方也已。」

《论语·颜渊篇第十二 12.16》 子曰：「君子成人之美，不成人之恶。小人反是。」

《论语·颜渊篇第十二 12.24》 曾子曰：「君子以文会友，以友辅仁。」

《论语·子路篇第十三 13.23》 子曰：「君子和而不同，小人同而不和。」

《论语·学而篇第一 1.16》 子曰：「不患人之不己知，患不知人也。」

《论语·宪问篇第十四 14.30》 子曰：「不患人之不己知，患其不能也。」
\end{classic}

\entrydate{2017.07.08}

\chapter{《关睢》与“兴”的表达手法}

\begin{classic}
“关关雎鸠，在河之洲”，毛传：“兴也。”所谓“兴”，可以说是引起话题，或者说是由景引起情。这景与情的碰合多半是诗人当下的感悟，它可以是即目，也不妨是浮想；前者是实景，后者则是心象。但它仅仅是引起话题，一旦进入话题，便可以放过一边，因此“兴”中并不含直接的比喻，若然，则即为“比”。——《诗经别裁》摘录
\end{classic}

诗经多用兴的手法，而之后古人说话也多引用诗经中的句子来“兴”起自己要表达的意思，从而使表达变得委婉文雅。比如，两国之间外交官的官方对话必要以诗经作为引子，而且还会奏乐——因为诗经原本都是可以唱的。古代人说话不会太直白的这一特点也值得今人借鉴。比如，孔子的学生端木赐想要知道孔子是否想从政，直接问太唐突，他便说木匣中有一美玉，请问是将其卖掉呢，还是自己珍藏呢？端木赐，名子贡，是孔子门下言语科排名第一的学生，口才极好，由此可见一斑。“兴”的表达方式也可见于文人将其志寄寓于植物。如松柏于岁寒中后凋，竹之坚韧正直。

对于《关睢》的理解也应借助“兴”升华到更高的层次与境界。即，将“窈窕淑女”引申为我们的终极理想、至善之地，而借君子“寤寐思服”、“辗转反侧”的“好逑”兴起自己对志向矢志不渝的追求。这样，一首诗便不仅仅是具有美感的艺术，更是对我们人生中的关键因素，如立志、践行、时穷节乃见、不轻易放弃，有着重要的现实推动意义。而前提是需要提升我们自身的领悟能力，使其能够与古人的优秀思想对接。

\entrydate{2017.01.21}

\chapter{回归文字的原义}

对于《论语》、《诗经》等作品所蕴含道理的理解仅停留于字面显然是不行的。但若拘泥于引申后正确、标准的现代含义也仍旧不足。比如，今晨醒来后脑海浮现出了《论语 学而第一 1.15》中的那句“如切如磋，如琢如磨”（源于《诗经 卫风 淇奥》）。由于对这句话已经耳熟能详，甚至都产生免疫力了，所以对于“切磋琢磨”四个字的第一感觉便是我们常用的引申义，即，通过不断打磨和完善自我，以追求君子的完美境界。这样的理解固然不错，但却丧失了文字原本鲜活的形象和生动的意义。因而，就像是只有法律条文的约束但缺乏宗教情怀的感召与启发一样，如此的理解对于指导实践、激励个人立志、促进自我提升的作用反倒不那么显著与深刻了。

周振甫《诗经译注》中对“切磋琢磨”的解释为：“治骨曰切，治象牙曰磋。治玉曰琢，治石曰磨”。若将这几件工作的具体形象呈现出来，就会发现将其用于形容君子择善固执的提升过程与止于至善的终极境界是多么地恰如其分了。可以想象，一位技艺卓越、一丝不苟的工匠，他那经过了深思熟虑后的每一刀切下去都是无比地稳定沉着、精准无误。骨角玉石在他每一次小心翼翼但又娴熟的雕琢与打磨后虽只落下那么一点点亮晶晶的粉屑，但如此经过成百上千次地反复后，外形、质地、纹理上细致入微的变化终于累积出了巧夺天工般的维妙维肖与圆润细腻。很显然，好习惯由个别到完整体系的逐步养成，以及最终君子品性的塑成，都不是一蹴而就的，都是通过这么一点点地精雕细琢，由无数个微小细节积累出来的。每一次的提升与改进虽然微小，但却无比扎实，毫无后顾之忧。这便是任何人想要臻于完美所应遵循的一般过程。任何有悖于此的浮躁、急切、贪心、自以为是，无论能够从中得到多少现实的效果与回报，都是没有意义且难以持久的。

最后，让我们再回顾一下《论语 学而第一 1.15》的原文和杨伯峻的翻译吧！

原文：

\begin{classic}
子贡曰：“诗云：‘如切如磋，如琢如磨’，其斯之谓与？”子曰：“赐也，始可与言诗已矣，告诸往而知来者。”
\end{classic}

翻译：

\begin{classic}
子贡说：“诗经上说：‘要像对待骨、角、象牙、玉石一样，先开料，再糙锉，细刻，然后磨光。’那就是这样的意思吧？”孔子道：“赐呀，现在可以同你讨论诗经了，告诉你一件，你能有所发挥，举一反三了。”

《论语·学而篇第一 1.15》 子贡曰：「贫而无谄，富而无骄，何如？」子曰：「可也，未若贫而乐，富而好礼者也。」子贡曰：「诗云：『如切如磋，如琢如磨』，其斯之谓与？」子曰：「赐也，始可与言诗已矣，告诸往而知来者。」
\end{classic}

\entrydate{2018.07.29}

\chapter{性相近，习相远}

\begin{classic}
《论语·阳货篇第十七 17.2》 子曰：「性相近也，习相远也。」
\end{classic}

性相近，是指人性大都具有欣赏、向往、追求真、善、美的倾向。其中，“真”代表科学与理性，“善”为道德与正义，“美”是艺术性的升华（exaltation）。这三者并非泾渭分明，而是各自均在一定程度上包含其他成分，相辅相成。这种人性的倾向并不一定最终达成事实上的必然，仅代表了人性的 \textbf{初值} 和今后发展的一种可能。因此，由于每个人后天学习、经验、磨炼、成长的路径不同，志向与努力的程度也不同，并由此导致言语、思想、行动、习惯、命运的差异，最终渐行渐远，鲜有殊途同归。这便是《论语·子罕篇第九 9.30》所讲的：

\begin{classic}
子曰：「可与共学，未可与适道；可与适道，未可与立；可与立，未可与权。」
\end{classic}

本章“习相远”中的“习”指人的习惯、习性。同时，我们也还记得，《论语·学而篇第一 1.1》中孔子讲“学而时习”，《论语·学而篇第一 1.4》中曾子讲“传不习乎”，其中的“习”均指练习、实践。若将其与本章的“习相远”对照，则可看出不同的人在后天学习与实践上花费的工夫与修炼的程度往往存在很大差异。这样一来，“三人行必有我师”，我们总能在某方面找到让自己学习、提升的榜样，也会发现让自己引以为戒的反面教材。因此，“习相远”这一客观事实，有助于我们形成“见贤思齐，见不贤而内自省”的动力。

\entrydate{2023.04.19}

\chapter{故作高深}

\begin{classic}
《论语·宪问篇第十四 14.35》子曰：「莫我知也夫！」子贡曰：「何为其莫知子也？」子曰：「不怨天，不尤人，下学而上达。知我者，其天乎！」
\end{classic}

虽然孔夫子与苏格拉底同属人类文明轴心时代的四大圣哲，其与各自弟子对话、探讨的方式却完全不同。前者故弄玄虚、故作高深、曲折逶迤。后者分类定义、逻辑严谨、平等透彻。由此，在探求真理、求知除昧的层面，他们对于人类文明的影响与推动效果迥异。

孔夫子与弟子们的谈话方式给后世直至今日的中国人在为政、治学、实业等诸多领域带来了负面影响。看似高深莫测、自我感觉真诚而良好，但往往实际说的是大话、空话、漂亮话（即邓晓芒老师所说的自欺）。虽然字句对仗工整，但内在逻辑链条实则很短，更无所谓复杂体系的构建；概念没有明确定义，相互重叠且整体上不完备。在教与学的具体实践上，老师的讲解、讲义、教材鲜有深入、透彻、酣畅淋漓者。学生中自由探讨，大胆质疑，求真理而非褒吾师者亦不多。

论语中类似打哑谜的对话还有：

\begin{classic}
《论语·八佾篇第三 3.13》 王孙贾问曰：「与其媚于奥，宁媚于灶，何谓也？」子曰：「不然，获罪于天，无所祷也。」

《论语·述而篇第七 7.15》 冉有曰：「夫子为卫君乎？」子贡曰：「诺，吾将问之。」入曰：「伯夷、叔齐何人也？」曰：「古之贤人也。」曰：「怨乎？」曰：「求仁而得仁，又何怨？」出曰：「夫子不为也。」
\end{classic}

\entrydate{2023.05.07}

\chapter{《论语》中的生动画面}

\begin{classic}
《论语·乡党篇第十 10.27》 色斯举矣，翔而后集。曰：「山梁雌雉，时哉时哉！」子路共之，三嗅而作。
\end{classic}

这一章描绘出了一幅生动的画面：深秋时节，孔子与弟子们在层林尽染的山间漫步。道旁是潺潺的溪流，溪水清冽。一群雌雉在桥边觅食，时而飞起散开，时而降落聚集。在这清冷、萧索的天地之间，孔子的头脑清醒、精神振奋，有感于大自然中的鸟兽生灵，获得了君子“时哉”的启发：君子应把握时机，顺应时机，无可无不可，但一切又应以坚守的原则和远大的终极理想为前提。难怪，孟子后来将孔子称作“圣之时者也”。闻过则喜的子路则在内心暗暗记下了老师的感悟，以示受教的谢意——这正是他的可爱与纯真之处。

\begin{center}
\includegraphics[width=\textwidth,height=0.62\textheight,keepaspectratio]{figures/fig03.png}
\end{center}

\entrydate{2023.09.18}

\chapter{克己复理}

\begin{classic}
《论语·颜渊篇第十二 12.1》颜渊问仁。子曰：「克己复礼，为仁。一日克己复礼，天下归仁焉。为仁由己，而由仁乎哉？」颜渊曰：「请问其目？」子曰：「非礼勿视，非礼勿听，非礼勿言，非礼勿动。」颜渊曰：「回虽不敏，请事斯语矣！」
\end{classic}

将此章中的“克己复礼”改为“克己复理”就对了！

“克己”的意思是“自己能够”，即调动自己的主观能动性，将天赋、潜力、勤奋发挥到极致。英文里对应的表述则是take your initiatives：

\begin{classic}
take your initiatives (initiative [Oxford Advanced Learner] n. the ability to decide and act on your own without waiting for sb to tell you what to do.)。
\end{classic}

“复理”与“复礼”的意思正好相反，是去除人与人之间的等级、贵贱差异——例如在学术领域就是要消除学阀与青年学者、大团队与小团体在项目申请与评审上的潜在不公平性——代之以在理性的指引下，基于具有普遍性与必然性的原则思考、分析、表达、做事。“使我们自己能够去做”而不是“在体制与政策的指引下去做”，从中体现出的就是自由、平等、去中心化、去权威化的精神。“独立之精神、自由之思想”，这是做学问的根本。

“克己”是从我们自身方面着眼，始终努力实践的基本原则。“复理”则是从外在环境与体制来说的，我们个人便很难对其施加影响以改变方向。环境好的时候，顺水行舟，邦有道则谷。环境不好的时候，仍应择善固执，避免成为体制与时代的牺牲品。无论外界如何变化，都应做出一些或者哪怕一点点能够真正发人深省、震撼人心、引领时代与推动人类的成果、产品、佳作。我想只有如此，才能无愧于清华“立德立言，无问西东”的精神。

\entrydate{2024.02.04}
